\capitulo{5}{Resultados}

\section{Resumen de resultados}
El trabajo realizado ha dado como resultado una herramienta web operativa que funciona como un asistente inteligente para la recomendación de Sistemas Aumentativos y Alternativos de Comunicación (SAAC). A través de un cuestionario sencillo, la aplicación es capaz de traducir las necesidades y capacidades de una persona en una propuesta concreta de soluciones tecnológicas.

\subsection{La aplicación web y el modelo de recomendación}
En lugar de obligar al usuario o al profesional a buscar entre decenas de catálogos, la herramienta centraliza una gran variedad de información de los SAAC del mercado en una base de datos. El diseño se ha centrado en que la interfaz sea muy limpia e intuitiva, pensada para usuarios no expertos.

\subsection{Pruebas con casos reales}
Para comprobar el funcionamiento real de la herramienta se introdujeron y analizaron los perfiles de dos pacientes reales evaluados en mayo de 2026 gracias a la colaboración de ELACyL. Los resultados de estas simulaciones muestran cómo responde el sistema ante dos casos diferentes:

\begin{itemize}
	\item \textbf{Paciente PAC-T5NE4D (Caso avanzado):}  Este perfil representa a una persona con un nivel de habla nulo (nivel 0) y fatiga pronunciada (resistencia nivel 1), aunque ve, oye y se maneja con la tecnología de forma aceptable (niveles 2). Como no tiene movilidad fina en 			las manos pero necesita comunicarse escribiendo (método alfabeto), el sistema detectó que la única forma de interactuar era a través de la mirada o movimientos de la cabeza, buscando programas compatibles con Windows y Android.
   	El resultado fue una lista muy completa. En cuanto a programas, recomendó herramientas potentes de comunicación como \textit{Grid 3, Communicator 5, Verbo, Look to Speak} y \textit{OptiKey}, además de \textit{Look to Learn} para entrenar el uso de la mirada. Para el equipo 			físico 	hardware), la aplicación seleccionó una tablet Android, un ordenador portátil, un lector de mirada (\textit{Eye tracker}), un joystick de mentón (\textit{BJOY Chin Plus}), varios pulsadores (de pedal y de soplido) y un brazo articulado para sujetarlo todo de forma segura en casa.
	\imagen{../img/caso_real_1.png}{Primer caso de paciente real.}{1}
    	La respuesta de feedback recibida por parte del paciente real validó el acierto del programa: el usuario aseguró que sí utiliza una de las combinaciones recomendadas y está ``muy satisfecho'' con esta. A su vez indica que está abierto a cambiar su sistema actual por otro en el 			futuro, por lo que confirma que el resto de soluciones propuestas le resultaron interesantes.
	
    	\item \textbf{Paciente PAC-7GUYVR (Caso de ELA bulbar):}  Este segundo perfil es el opuesto al anterior. El paciente presenta un caso de ELA de inicio bulbar, lo cual significa que el habla se pierde antes que la movilidad. Ve y oye perfectamente (niveles 3), aún conserva algo de 			habla (nivel 2), tiene mucha resistencia (nivel 3) y se mueve en silla de ruedas, pero tiene total autonomía en sus decisiones (nivel 3). Como puede usar las manos para interactuar, el sistema simplificó la recomendación al máximo hacia una aplicación de escritura a voz.
    	Por el lado del software, solo recomendó un programa básico de \textit{Asistente de voz AAC}. Por el lado del hardware, sugirió: un \textit{iPad}, una tablet Android, un portátil (ya que el usuario no seleccionó una plataforma de preferencia y esta aplicación se encuentra disponible  		multiplataforma), un ratón Bluetooth y un lápiz táctil, añadiendo un soporte para anclar la pantalla a la silla de ruedas (ya que marcó que sí utilizaba).
	\imagen{../img/caso_real_2.png}{Segundo caso de paciente real.}{1}
    	Aunque este paciente también definió la experiencia como ``muy satisfecho'', marcó que ``No'' utiliza la recomendación de forma habitual. La posible razón tras de esto sea que, al coservar un buen nivel de habla e independencia, ve el SAAC como un apoyo extra o 					puntual, y no como una necesidad vital para su día a día, pero que podrá considerar cuando su enfermedad progrese.
\end{itemize}
Estas dos pruebas demuestran que la herramienta web funciona correctamente, no da respuestas genéricas, sino que sabe cuándo proponer un despliegue tecnológico más grande y cuándo sugerir una solución sencilla y directa.

\subsection{Pruebas con casos simulados}
Debido al reducido tamaño del grupo de estudio se debieron realizar también pruebas de simulación para comprobar el funcionamiento en otra diversidad de situaciones. A continuación se mostrarán las más relevantes y aquellas con parámetros más diferenciados de los anteriores:

\begin{itemize}
	    \item \textbf{Caso de prueba 1: PAC-XPIOTW (Acceso por voz):}  Simulamos el perfil de un paciente en las primeras etapas de la enfermedad. Mantiene sus capacidades sensoriales, tecnológicas y cognitivas al máximo (niveles 3) y total autonomía (nivel 3), pero su resistencia 			física empieza a flaquear (nivel 2). Como todavía habla perfectamente (nivel 3), el sistema aprovechó esta ventaja y se configuró para un entorno mixto mediante el acceso por voz, sin importar la plataforma.
    	El resultado orientó las soluciones hacia la autonomía del paciente en casa sin esfuerzo físico. En cuanto a software, recomendó sistemas de control por voz como \textit{Voice Access} y asistentes virtuales (\textit{Amazon Alexa / Google Home}). Para el apartado físico (hardware), el 		sistema sugirió una tablet Android montada en un soporte de anclaje articulado y un \textit{Enchufe Inteligente WiFi} para poder encender y apagar aparatos de la casa directamente con la voz.
	\imagen{../img/caso_prueba_1.png}{Primer caso de paciente simulado.}{1}

	 \item \textbf{Caso de prueba 2: PAC-JOUN1Y (Baja tecnología):}  Este caso se diseñó para simular una fase más avanzada donde el paciente no puede hablar en absoluto (nivel 0), su resistencia es moderada (nivel 2) y no tiene conocimientos tecnológicos o prefiere no usar 				pantallas (tecnología 0). Al seleccionar la plataforma Panel físico / Papel en un entorno domiciliario, el algoritmo descartó por completo los ordenadores y los programas informáticos recomendando únicamente paneles y tableros físicos.
    	Como soporte de hardware, la web prescribió herramientas de señalización física de gran utilidad: un \textit{Puntero Láser}, unas \textit{Gafas con Puntero Láser} integrado y un soporte de anclaje articulado para fijar los tableros.
	\imagen{../img/caso_prueba_2.png}{Segundo caso de paciente simulado.}{1}

	\item \textbf{Caso de prueba 3: PAC-SKNTHB (Validación lingüística):}  Este caso se diseñó con un doble objetivo: comprobar cómo se comporta el sistema cuando el usuario no define un método o interacción preferidos y, además, evaluar el soporte lingüístico seleccionando el 			\textit{Euskera} como idioma principal. El perfil simula a un paciente con capacidades intermedias (visión, resistencia y tecnología en nivel 2) y una buena autonomía general (nivel 3), que todavía conserva algo de habla (nivel 2), pero que busca opciones abiertas compatibles con 			cualquier plataforma en un entorno mixto.
	Como el sistema detectó que el paciente aún tiene habla residual y movilidad intermedia, ofreció un catálogo de software compatible con el euskera recomendando sistemas clave como \textit{Verbo, TD Snap} y \textit{Boardmaker 7}, junto al comunicador						\textit{QuickTalker 23}.
    	En cuanto al hardware, al no haber restricciones en la interacción, la herramienta desplegó una amplia gama de opciones para que el profesional o el usuario puedan elegir la que mejor se adapte a su evolución: desde dispositivos comunes (\textit{iPad}, tablet Android, ordenador 		portátil y ratón Bluetooth) hasta periféricos de acceso avanzado por si la enfermedad progresa, como un \textit{Eye tracker}, el joystick de mentón \textit{BJOY Chin Plus}, un soporte articulado y varios tipos de conmutadores (el \textit{Spec Amarillo}, de pedal y de soplido).
	\imagen{../img/caso_prueba_3.png}{Tercer caso de paciente simulado.}{1}

	\item \textbf{Caso de prueba 4: PAC-L1O0R7 (Sin resultados compatibles):} Este caso se diseñó con un fin de control crítico: evaluar el algoritmo ante un perfil clínico con capacidades extremadamente limitadas y criterios excluyentes. Se simuló a un  paciente en una fase de máxima 	afectación orgánica y cognitiva, con niveles nulos en visión, audición, tecnología y habla (todos en nivel 0) y una fatiga física crítica (nivel 0). 
    	Como era de esperar debido a las estrictas restricciones relacionales de la base de datos, el sistema no devolvió ninguna coincidencia, mostrando una pantalla de advertencia sin software, hardware ni servicios disponibles. .
	\imagen{../img/caso_prueba_4.png}{Cuarto caso de paciente simulado.}{1}
\end{itemize}